\chapter{État de l'art}

Avant de définir les fonctions attendues de l'appareil, ce chapitre situe le projet par rapport à ce qui existe déjà. Il retrace d'abord l'évolution du lecteur audio portable. Il expose ensuite les solutions techniques disponibles, puis le positionnement du présent travail.


\section{Travaux existants}

\subsection{Du baladeur au lecteur audio dédié}
Le lecteur audio portable s'impose dès la fin des années 1970 avec le baladeur à cassette. Le disque compact portable lui succède, puis la lecture de fichiers compressés bouleverse l'usage au tournant des années 2000 \cite{Wikipedia_Portable_Media_Player}. Le smartphone absorbe ensuite cette catégorie. Il réunit dans un seul appareil la lecture audio, la connectivité et le stockage. Le lecteur dédié a ainsi quitté le marché grand public \cite{Tom-Smigla_20_Best_DAP}.

\subsection{Le marché actuel, deux extrémités}
Le produit s'est recentré sur deux extrémités, séparées par un large écart de prix.

D'un côté subsiste un segment audiophile actif, porté par des marques comme Astell\&Kern, FiiO, HiBy ou iBasso. Ces appareils, appelés \emph{DAP} (Digital Audio Player), visent la très haute fidélité, avec plusieurs convertisseurs en parallèle et des sorties symétriques, pour quelques centaines à plusieurs milliers de francs \cite{Tom-Smigla_20_Best_DAP, Homecine_DAP_Comparison}. Beaucoup exécutent Android sur un processeur applicatif et se rapprochent d'un smartphone mais sans fonction téléphonique.

De l'autre se maintiennent des lecteurs d'entrée de gamme, simples, peu coûteux et peu gourmands, destinés à un usage quotidien sans exigence audiophile \cite{AudiophileOn_Best_Cheap_DAP}.


\section{Technologies disponibles}
Cette section expose les options techniques offertes pour les grandes fonctions d'un tel appareil. Elle ne tranche pas. Les choix retenus sont établis plus loin, une fois le cahier des charges posé.

\subsection{Architectures de calcul}
Deux approches coexistent. Un processeur applicatif exécutant un système d'exploitation complet apporte de la souplesse et l'accès au streaming, au prix d'une consommation et d'une complexité élevées \cite{Homecine_DAP_Comparison}. Un microcontrôleur suffit pour un appareil à fonction unique. Il consomme peu et son comportement temporel reste maîtrisé, mais le logiciel se développe au plus près du matériel.

Le partage entre les deux se fait sur une question simple. L'appareil doit-il accéder à des services en ligne, ou seulement lire des fichiers locaux ?

\subsection{Conversion et amplification}
La chaîne de restitution décode le flux, le convertit en analogique, le filtre puis l'amplifie. Deux décisions la structurent.

La première porte sur la sortie du convertisseur. Une sortie asymétrique transmet le signal sur un seul conducteur et demande peu de composants. Une sortie différentielle le transmet sur deux conducteurs opposés. L'étage suivant ne retient que leur différence, ce qui rejette les perturbations communes aux deux lignes et double l'amplitude utile, au prix de composants supplémentaires.

La seconde porte sur le découpage de la chaîne. Un circuit unique réunissant le convertisseur et l'amplificateur réduit le nombre de composants, mais supprime tout point de mesure entre les deux étages.

\subsection{Transport sans fil}
La sortie filaire par jack préserve le signal analogique. La liaison Bluetooth est plus pratique, mais elle subit la compression et les limites de la portée radio. Le profil A2DP y transporte le flux stéréo et l'AVRCP assure la télécommande \cite{bluetooth_sig_a2dp}.

Le choix du codec est le point sensible. Le SBC est le seul dont la prise en charge est obligatoire, ce qui garantit qu'il fonctionne avec tout récepteur \cite{bluetooth_sig_a2dp}. L'AAC, l'aptX et le LDAC offrent un meilleur compromis entre débit et qualité, mais ils sont propriétaires ou soumis à licence et ne sont pas implémentés partout.

\subsection{Énergie et stockage}
L'alimentation repose sur un accumulateur lithium-ion, un circuit de charge et, dans les conceptions soignées, une jauge qui estime l'état de la batterie. La recharge passe aujourd'hui majoritairement par USB-C.

Le stockage suit deux voies. Une mémoire interne fixe la capacité à la fabrication et impose un câble pour changer la bibliothèque. Une carte microSD amovible lève ces deux limites, mais demande un connecteur accessible et un système de fichiers standard.


\section{Positionnement du projet}
Le présent travail ne relève d'aucune des deux catégories commerciales décrites plus haut. Il ne vise ni les performances des DAP audiophiles, ni la polyvalence d'un smartphone.

Son objectif est de concevoir un appareil dédié, compact et fonctionnel, de bout en bout et à des fins formatrices. Chaque choix matériel et logiciel doit être maîtrisé et justifié. La contribution recherchée n'est pas un produit inédit sur le marché, mais la démarche complète d'intégration, du composant jusqu'au firmware.


\begin{table}[h]
      \centering
      \renewcommand{\arraystretch}{1.15}
      \setlength{\tabcolsep}{8pt}
      \footnotesize
      \begin{tabular}{m{2.6cm} | m{3.4cm} | m{3.4cm} | m{3.4cm}}
            \hline
            \textbf{Critère} & \textbf{DAP audiophile}         & \textbf{DAP entrée de gamme} & \textbf{Ce projet}                   \\
            \hline
            Calcul           & Processeur applicatif (Android) & Microcontrôleur              & Microcontrôleur                      \\
            Objectif audio   & Très haute fidélité             & Basique                      & Performances mesurées et documentées \\
            Sorties          & Jack et Bluetooth               & Jack et Bluetooth            & Jack et Bluetooth                    \\
            Finalité         & Produit commercial              & Produit commercial           & Prototype académique                 \\
            \hline
      \end{tabular}
      \caption{Positionnement du projet par rapport aux solutions existantes.}
      \label{tab:comparatif}
\end{table}

Les fonctions attendues de l'appareil découlent de ce positionnement. Elles font l'objet du chapitre suivant.